\documentclass[oneside,a4paper,14pt]{extarticle}
\usepackage[T1,T2A]{fontenc}
\usepackage[a4paper,
letterpaper,top=20mm,bottom=20mm,
left=30mm,right=10mm]{geometry}
\usepackage[utf8]{inputenc}
\usepackage[russian]{babel}
\usepackage{textcomp}
\usepackage{indentfirst}
\usepackage{graphicx}
\usepackage{mwe}
\usepackage{pdfpages}
\usepackage{array}
\usepackage{booktabs}
\usepackage{adjustbox}
\usepackage{caption}
\usepackage{comment}
\usepackage{listings}
\usepackage{float}
\usepackage{xcolor}
\usepackage{hyperref}
\usepackage{titlesec}

\captionsetup[figure]{name = {Рисунок}, labelsep=endash}
\titleformat{\section}{\normalsize\bfseries}{\thesection}{1em}{}
\titleformat{\subsection}{\normalsize\bfseries}{\thesubsection}{1em}{}

\renewcommand\baselinestretch{1.33}\normalsize
\setlength{\parindent}{1.25cm}

\lstset{
    frame=tb,
    extendedchars=false,
    inputencoding=utf8,
    numbers=left,
    language=Python,
    aboveskip=3mm,
    belowskip=3mm,
    showstringspaces=false,
    basicstyle=\small\ttfamily,
    breaklines=true,
    breakatwhitespace=false,
    columns=fullflexible,
    tabsize=2,
    keepspaces=true,
    linewidth=\linewidth,
    keywordstyle=\color{blue},
    commentstyle=\color{gray},
    stringstyle=\color{orange},
}

\begin{document}
\pagestyle{empty}
\begin{center}
    {МИНИСТЕРСТВО НАУКИ И ВЫСШЕГО ОБРАЗОВАНИЯ РОССИЙСКОЙ ФЕДЕРАЦИИ \\
    Федеральное государственное бюджетное образовательное учреждение высшего образования \\
    «Вятский государственный университет» (ФГБОУ ВО «ВятГУ») \\
    Факультет автоматики и вычислительной техники \\
    Кафедра электронных вычислительных машин \\ }

    \vspace*{35mm}

    {«Сценарное и функциональное программирование» \\
    Лабораторная работа №2 \\ }
\end{center}

\vspace*{35mm}
       Выполнил студент группы ИВТб-3302-05-00 \hspace*{1mm} \underline{\hspace{10mm}} /Фофанов А.С./  \\
       \hspace*{12mm}Проверил преподаватель \hspace*{44mm} \underline{\hspace{10mm}} /\hspace{30mm}/ \\

\begin{center}
    \vspace*{35mm}
    \large{Киров \\
    2026}
\end{center}

\newpage
\pagestyle{plain}
\setcounter{page}{2}
\normalsize

\section*{Цель работы}
Разработать костяк backend-сервиса на языке Python для управления пользователями и OAuth-клиентами. Сервис должен предоставлять GraphQL API, реализовывать аутентификацию по JWT (HS256), безопасно хранить пароли с использованием Argon2id и pepper, вести аудит действий пользователя и предоставлять минимальную админ-панель.

\section*{Задание}
\begin{enumerate}
    \item Спроектировать слой данных на SQLAlchemy 2.x и настроить миграции через Alembic.
    \item Организовать FastAPI-приложение с разделением на модули.
    \item Реализовать GraphQL API на библиотеке Ariadne (schema-first: SDL-схема, резолверы, контекст, директивы-пермишены).
    \item Реализовать базовую аутентификацию по JWT (HS256) без OAuth2.
    \item Реализовать безопасное хэширование паролей: Argon2id, глобальный pepper, оценка стойкости через zxcvbn.
    \item Реализовать аудит действий пользователя в таблице audit\_events.
    \item Создать минимальный SPA-интерфейс администратора на vanilla JS.
    \item Покрыть ключевые сценарии тестами на pytest.
\end{enumerate}

\section*{Описание реализации}

Проект состоит из следующих модулей:

\begin{itemize}
    \item \texttt{main.py} — точка входа FastAPI-приложения, подключает роутеры и монтирует ASGI-приложения;
    \item \texttt{config.py} — конфигурация через pydantic-settings, читает \texttt{.env};
    \item \texttt{db/base.py}, \texttt{db/session.py} — декларативный Base SQLAlchemy, движок, фабрика сессий, зависимость \texttt{get\_db};
    \item \texttt{models/user.py}, \texttt{models/oauth\_client.py}, \texttt{models/audit\_event.py} — ORM-модели;
    \item \texttt{core/passwords.py} — хэширование Argon2id, pepper, zxcvbn;
    \item \texttt{core/jwt.py} — выпуск и проверка JWT (HS256);
    \item \texttt{core/deps.py} — зависимости FastAPI для аутентификации;
    \item \texttt{auth/router.py}, \texttt{auth/service.py}, \texttt{auth/schemas.py} — REST-эндпоинты /auth/*;
    \item \texttt{audit/service.py}, \texttt{audit/events.py} — единая точка записи событий аудита;
    \item \texttt{graphql/schema.graphql}, \texttt{schema.py}, \texttt{resolvers.py}, \texttt{context.py}, \texttt{directives.py} — GraphQL API на Ariadne;
    \item \texttt{static/index.html}, \texttt{app.js}, \texttt{style.css} — мини-админка на vanilla JS.
\end{itemize}

Конфигурация приложения хранится в файле \texttt{.env} и включает: строку подключения к базе данных, секрет для подписи JWT, алгоритм, время жизни токена, pepper для паролей, минимальную длину пароля.

\subsection*{Модуль core/passwords.py}

Модуль реализует хэширование паролей с использованием Argon2id. Перед хэшированием к паролю добавляется глобальный pepper, хранящийся в переменных окружения. Оценка стойкости пароля выполняется через zxcvbn с требованием score не ниже 3.

\begin{lstlisting}
_hasher = PasswordHasher(
    time_cost=3,
    memory_cost=65536,
    parallelism=4,
    hash_len=32,
    salt_len=16,
)

def hash_password(password: str) -> str:
    if not isinstance(password, str) or not password:
        raise WeakPasswordError("Пароль не может быть пустым.")
    peppered = password + settings.password_pepper
    return _hasher.hash(peppered)

def validate_policy(password, user_inputs=None):
    if len(password) < settings.password_min_length:
        raise WeakPasswordError("Пароль не должен быть короче N символов.")
    result = zxcvbn(password, user_inputs=user_inputs or [])
    if result["score"] < 3:
        raise WeakPasswordError("Слишком слабый пароль.")

\subsection*{Модуль core/jwt.py}

Модуль выпускает и проверяет JWT-токены с использованием HS256. В payload включаются claims \texttt{sub}, \texttt{iat}, \texttt{exp}, \texttt{jti}, \texttt{typ}. Поле \texttt{sub} приводится к строке в соответствии с RFC 7519. При декодировании проверяется подпись и срок действия.

\begin{lstlisting}
def create_access_token(sub, ttl_min=None):
    if ttl_min is None:
        ttl_min = settings.access_token_ttl_min
    now = datetime.now(timezone.utc)
    payload = {
        "sub": str(sub),
        "iat": int(now.timestamp()),
        "exp": int((now + timedelta(minutes=ttl_min)).timestamp()),
        "jti": uuid4().hex,
        "typ": "access",
    }
    return jwt.encode(payload, settings.jwt_secret,
                      algorithm=settings.jwt_algorithm)

def decode_token(token: str) -> dict:
    if not token or not isinstance(token, str):
        raise ValueError("Пустой токен.")
    try:
        payload = jwt.decode(token, settings.jwt_secret,
                             algorithms=[settings.jwt_algorithm])
    except JWTError as e:
        raise ValueError("Невалидный токен") from e
    if payload.get("typ") != "access":
        raise ValueError("Не access-токен.")
    return payload
\end{lstlisting}

\subsection*{Модуль auth/service.py}

Сервисный слой содержит бизнес-логику регистрации и логина, используется как REST-роутером, так и GraphQL-резолвером. При неудачном логине всегда возвращается одно и то же исключение \texttt{InvalidCredentialsError}, что не позволяет клиенту различить случаи «несуществующий email» и «неверный пароль».

\begin{lstlisting}
def register_user(db, *, email, password, request=None):
    email_norm = email.strip().lower()
    existing = db.execute(
        select(User).where(User.email == email_norm)
    ).scalar_one_or_none()
    if existing is not None:
        raise EmailAlreadyExistsError("Email already registered")
    validate_policy(password, user_inputs=[email_norm])
    user = User(
        email=email_norm,
        hashed_password=hash_password(password),
        is_active=True, is_admin=False,
    )
    db.add(user)
    db.flush()
    write_event(db, UserEvents.CREATED, actor_id=user.id, ...)
    token = create_access_token(sub=user.id)
    return user, token
\end{lstlisting}

\subsection*{Модели базы данных}

Модель \texttt{User} содержит email (уникальный, индексируемый), хэш пароля, флаги \texttt{is\_active} и \texttt{is\_admin}, дату создания. Модель \texttt{OAuthClient} хранит метаданные клиента: \texttt{client\_id}, имя, флаг конфиденциальности, список \texttt{redirect\_uris}. Модель \texttt{AuditEvent} фиксирует события: тип, актора, IP, User-Agent, произвольные metadata в формате JSON.

\begin{lstlisting}
class User(Base):
    __tablename__ = "users"
    id: Mapped[int] = mapped_column(primary_key=True, autoincrement=True)
    email: Mapped[str] = mapped_column(String(255), unique=True, index=True)
    hashed_password: Mapped[str] = mapped_column(String(255))
    is_active: Mapped[bool] = mapped_column(Boolean, default=True)
    is_admin: Mapped[bool] = mapped_column(Boolean, default=False)
    created_at: Mapped[datetime] = mapped_column(DateTime, default=_utcnow)

\subsection*{Модуль graphql/schema.py}

GraphQL API реализован на библиотеке Ariadne в schema-first стиле. Схема описана в отдельном файле \texttt{schema.graphql}. К полям Query и Mutation привязаны резолверы, реализованы директивы-пермишены \texttt{@isAuthenticated} и \texttt{@isAdmin}. Для преобразования имён полей из camelCase в snake\_case используется \texttt{snake\_case\_fallback\_resolvers}.

\begin{lstlisting}
schema = make_executable_schema(
    type_defs,
    query,
    mutation,
    audit_event_type,
    snake_case_fallback_resolvers,
    directives={
        "isAuthenticated": IsAuthenticatedDirective,
        "isAdmin": IsAdminDirective,
    },
)

graphql_app = GraphQL(
    schema,
    context_value=get_context_value,
    debug=True,
)
\end{lstlisting}

\subsection*{Модуль audit/service.py}

Аудит реализован через единую функцию \texttt{write\_event}, которая создаёт запись в таблице \texttt{audit\_events} и коммитит транзакцию. Имена событий соответствуют соглашению \texttt{<объект>.<действие>}. Для защиты от записи секретов в metadata предусмотрена проверка по списку запрещённых ключей.

\begin{lstlisting}
FORBIDDEN_METADATA_KEYS = frozenset({
    "password", "hashed_password", "token",
    "access_token", "refresh_token", "secret", "pepper", "jwt",
})

def write_event(db, event_type, *, actor_id=None, client_id=None,
                ip=None, user_agent=None, metadata=None):
    if metadata:
        forbidden = FORBIDDEN_METADATA_KEYS & metadata.keys()
        if forbidden:
            raise ValueError(f"Аудит не должен содержать секреты: {forbidden}")
    event = AuditEvent(
        event_type=event_type, actor_id=actor_id,
        client_id=client_id, ip=ip, user_agent=user_agent,
        metadata_json=metadata,
    )
    db.add(event)
    db.commit()
    db.refresh(event)
    return event
\end{lstlisting}

\subsection*{Мини-админка}

Админ-панель реализована на vanilla JS. Она взаимодействует с REST-эндпоинтом \texttt{/auth/login} для получения JWT и с GraphQL-эндпоинтом \texttt{/graphql/} для запросов к данным. Токен хранится в localStorage. В интерфейсе доступны: список OAuth-клиентов, создание клиента, просмотр профиля, список пользователей и отключение пользователя.

\begin{lstlisting}
async function gql(query, variables = {}) {
    const headers = { 'Content-Type': 'application/json' };
    const token = localStorage.getItem(TOKEN_KEY);
    if (token) headers['Authorization'] = 'Bearer ' + token;
    const response = await fetch('/graphql/', {
        method: 'POST', headers,
        body: JSON.stringify({ query, variables }),
    });
    return response.json();
}
\end{lstlisting}

\section*{Тестирование}

Ключевые сценарии покрыты тестами на pytest. Тесты используют отдельную базу данных SQLite, которая создаётся и удаляется автоматически. Фикстуры pytest обеспечивают изоляцию между тестами. Всего реализовано 42 теста, покрывающих: политику паролей, хэширование и проверку, REST-эндпоинты /auth/*, GraphQL-запросы и мутации, работу директив-пермишенов, ведение аудита.

\section*{Экранные формы}

\begin{figure}[H]
    \centering
    \includegraphics[width=1\linewidth]{1.png}
    \caption{Страница входа в админ-панель}
    \label{fig:admin-login}
\end{figure}

\begin{figure}[H]
    \centering
    \includegraphics[width=1\linewidth]{2.png}
    \caption{Таблица пользователей в админ-панели}
    \label{fig:admin-users}
\end{figure}

\begin{figure}[H]
    \centering
    \includegraphics[width=1\linewidth]{3.png}
    \caption{Swagger UI с REST-эндпоинтами}
    \label{fig:swagger}
\end{figure}

\clearpage
\section*{Вывод}
В ходе выполнения лабораторной работы был разработан костяк backend-сервиса Vault для управления пользователями и OAuth-клиентами. Реализованы: слой данных на SQLAlchemy 2.x с миграциями Alembic, REST-эндпоинты аутентификации на FastAPI, GraphQL API на Ariadne с директивами-пермишенами, безопасное хэширование паролей (Argon2id + pepper + zxcvbn), выпуск и проверка JWT (HS256), аудит значимых действий пользователя, минимальная админ-панель на vanilla JS. Все ключевые сценарии покрыты тестами. Сервис готов к расширению OAuth2/OIDC в следующей лабораторной работе.

\end{document}
